\documentclass[10pt,a4paper]{amsart}
\usepackage{amsmath,amssymb,amsthm,graphicx}
\usepackage{verbatim}
\newenvironment{reference}{\comment}{\endcomment}
%\newenvironment{reference}{}{}
\newenvironment{slogan}{\comment}{\endcomment}
\newenvironment{history}{\comment}{\endcomment}
% For commutative diagrams we use Xy-pic
\usepackage[all]{xy}

% We use 2cell for 2-commutative diagrams.
\xyoption{2cell}
% Theorem environments.
%
\theoremstyle{plain}
\newtheorem{theorem}[subsection]{Theorem}
\newtheorem{proposition}[subsection]{Proposition}
\newtheorem{lemma}[subsection]{Lemma}

\theoremstyle{definition}
\newtheorem{definition}[subsection]{Definition}
\newtheorem{example}[subsection]{Example}
\newtheorem{exercise}[subsection]{Exercise}
\newtheorem{situation}[subsection]{Situation}

\theoremstyle{remark}
\newtheorem{remark}[subsection]{Remark}
\newtheorem{remarks}[subsection]{Remarks}

\numberwithin{equation}{subsection}

% Macros
%
\def\lim{\mathop{\mathrm{lim}}\nolimits}
\def\colim{\mathop{\mathrm{colim}}\nolimits}
\def\Spec{\mathop{\mathrm{Spec}}}
\def\Hom{\mathop{\mathrm{Hom}}\nolimits}
\def\Ext{\mathop{\mathrm{Ext}}\nolimits}
\def\SheafHom{\mathop{\mathcal{H}\!\mathit{om}}\nolimits}
\def\SheafExt{\mathop{\mathcal{E}\!\mathit{xt}}\nolimits}
\def\Sch{\mathit{Sch}}
\def\Mor{\mathop{\mathrm{Mor}}\nolimits}
\def\Ob{\mathop{\mathrm{Ob}}\nolimits}
\def\Sh{\mathop{\mathit{Sh}}\nolimits}
\def\NL{\mathop{N\!L}\nolimits}
\def\CH{\mathop{\mathrm{CH}}\nolimits}
\def\proetale{{pro\text{-}\acute{e}tale}}
\def\etale{{\acute{e}tale}}
\def\QCoh{\mathit{QCoh}}
\def\Ker{\mathop{\mathrm{Ker}}}
\def\Im{\mathop{\mathrm{Im}}}
\def\Coker{\mathop{\mathrm{Coker}}}
\def\Coim{\mathop{\mathrm{Coim}}}

% Boxtimes
%
\DeclareMathSymbol{\boxtimes}{\mathbin}{AMSa}{"02}

%
% Macros for moduli stacks/spaces
%
\def\QCohstack{\mathcal{QC}\!\mathit{oh}}
\def\Cohstack{\mathcal{C}\!\mathit{oh}}
\def\Spacesstack{\mathcal{S}\!\mathit{paces}}
\def\Quotfunctor{\mathrm{Quot}}
\def\Hilbfunctor{\mathrm{Hilb}}
\def\Curvesstack{\mathcal{C}\!\mathit{urves}}
\def\Polarizedstack{\mathcal{P}\!\mathit{olarized}}
\def\Complexesstack{\mathcal{C}\!\mathit{omplexes}}
% \Pic is the operator that assigns to X its picard group, usage \Pic(X)
% \Picardstack_{X/B} denotes the Picard stack of X over B
% \Picardfunctor_{X/B} denotes the Picard functor of X over B
\def\Pic{\mathop{\mathrm{Pic}}\nolimits}
\def\Picardstack{\mathcal{P}\!\mathit{ic}}
\def\Picardfunctor{\mathrm{Pic}}
\def\Deformationcategory{\mathcal{D}\!\mathit{ef}}
\usepackage{cleveref}
\crefname{originallemma}{original lemma}{original lemmas}
\Crefname{originallemma}{Original lemma}{Original lemmas}
\makeatletter
\newlabel{lemma-chern-classes-dual}{{0FA5}{source}{}{sourcepointer-1}{}}
\newlabel{lemma-chern-classes-dual@cref}{{[originallemma][0][]0FA5}{[1][0][]source}{}{}{}}
\makeatother
\crefname{originallemma}{original lemma}{original lemmas}
\Crefname{originallemma}{Original lemma}{Original lemmas}
\makeatletter
\newlabel{lemma-commutative-chern-perfect}{{0FAC}{source}{}{sourcepointer-2}{}}
\newlabel{lemma-commutative-chern-perfect@cref}{{[originallemma][0][]0FAC}{[1][0][]source}{}{}{}}
\makeatother
\crefname{originalremark}{original remark}{original remarks}
\Crefname{originalremark}{Original remark}{Original remarks}
\makeatletter
\newlabel{remark-splitting-principle-perfect}{{0FAD}{source}{}{sourcepointer-3}{}}
\newlabel{remark-splitting-principle-perfect@cref}{{[originalremark][0][]0FAD}{[1][0][]source}{}{}{}}
\makeatother
\crefname{originalreference}{original reference}{original references}
\Crefname{originalreference}{Original reference}{Original references}
\makeatletter
\newlabel{situation-setup}{{02QL}{source}{}{sourcepointer-4}{}}
\newlabel{situation-setup@cref}{{[originalreference][0][]02QL}{[1][0][]source}{}{}{}}
\makeatother
\title{The Chern classes, the Chern character components and the power sums of a perfect object are multiplied by the sign \$(-1)\textasciicircum{}i\$ when one passes to the dual}
\author{The Stacks Project Authors}
\date{stacks-project:chow.tex:local-snapshot-20261012; GNU FDL 1.2 or later}
\begin{document}
\maketitle
\noindent\textit{Editorial note.} All delivered mathematical text is the author's own, copied from the pinned Stacks Project chapter file chow.tex; no statement or proof step was rewritten, reordered or removed.\par\medskip
\noindent\textit{Editorial note.} Preamble provenance: the pinned chapter file opens with \textbackslash{}input\{preamble\} and carries no preamble text of its own, so the project's single Stacks-wide preamble staged at sources/stacks-chow/preamble.tex is excerpted verbatim as three spans, lines 8-12 (the verbatim package together with the project's own reference, slogan and history display-environment definitions, including the author's commented-out older form of the reference environment), lines 14-18 (Xy-pic setup) and lines 154-217 (the native theorem declarations and mathematical macros). The lines 8-12 span is required because a counted statement or its surrounding text may use those display environments, which the project typesets as comments. Nothing was added, removed or redefined; the wrapper class and the amsmath/amssymb/amsthm/graphicx base packages are supplied by the assembler.\par\medskip
\noindent\textit{Editorial note.} Same-chapter pointers are located at their \textbackslash{}label lines inside the pinned chapter file; cross-chapter pointers are located in the project's own tag list staged in the snapshot (sources/stacks-chow/tags.txt), at the verbatim lines carrying their canonical tags.\par\medskip
\noindent\textit{Allowed standard prerequisites (editorial).}\begin{itemize}
\item The reduction of statements about a perfect object to a bounded complex of finite locally free modules of fixed rank. Same chapter chow.tex, label lemma-commutative-chern-perfect (invoked at proof line 9554). Its proof shows how to reduce to the case of a bounded complex of finite locally free modules of fixed rank, where the dual computation is available.
\item The dual statement for a finite locally free module. Same chapter chow.tex, label lemma-chern-classes-dual (invoked at proof line 9557). It supplies the sign \$(-1)\textasciicircum{}i\$ for the Chern classes of the dual in the reduced case.
\item The splitting principle for perfect objects, via Chern roots. Same chapter chow.tex, Remark \textbackslash{}ref\{remark-splitting-principle-perfect\} (invoked at proof line 9559). It gives the alternative route: the Chern roots of \$E\textasciicircum{}\textbackslash{}vee\$ are the negatives of those of \$E\$, which produces the same three signs.
\end{itemize}\medskip
\ifdefined\hypertarget\hypertarget{sourcepointer-1}{}\fi
\noindent\textit{Source reference.} Original lemma 0FA5: sources/stacks-chow/pinned-original-source.tex, source lines 8520--8520. Navigation pointer for \textbackslash{}ref\{lemma-chern-classes-dual\} used inside the counted span (line 9557); the argument that consumes it is retained in the span.\par\medskip
\ifdefined\hypertarget\hypertarget{sourcepointer-2}{}\fi
\noindent\textit{Source reference.} Original lemma 0FAC: sources/stacks-chow/pinned-original-source.tex, source lines 9394--9394. Navigation pointer for \textbackslash{}ref\{lemma-commutative-chern-perfect\} used inside the counted span (line 9554); the argument that consumes it is retained in the span.\par\medskip
\ifdefined\hypertarget\hypertarget{sourcepointer-3}{}\fi
\noindent\textit{Source reference.} Original remark 0FAD: sources/stacks-chow/pinned-original-source.tex, source lines 9494--9494. Navigation pointer for \textbackslash{}ref\{remark-splitting-principle-perfect\} used inside the counted span (line 9559); the argument that consumes it is retained in the span.\par\medskip
\ifdefined\hypertarget\hypertarget{sourcepointer-4}{}\fi
\noindent\textit{Source reference.} Original reference 02QL: sources/stacks-chow/pinned-original-source.tex, source lines 1372--1372. Navigation pointer for \textbackslash{}ref\{situation-setup\} used inside the counted span (line 9546); the argument that consumes it is retained in the span.\par\medskip
\begin{lemma}
\label{lemma-chern-classes-perfect-dual}
In Situation \ref{situation-setup} let $X$ be locally of finite type over $S$.
Let $E \in D(\mathcal{O}_X)$ be a perfect object whose Chern classes are
defined. Then $c_i(E^\vee) = (-1)^i c_i(E)$, $P_i(E^\vee) = (-1)^iP_i(E)$,
and $ch_i(E^\vee) = (-1)^ich_i(E)$ in $A^i(X)$.
\end{lemma}

\begin{proof}
First proof: argue as in the proof of
Lemma \ref{lemma-commutative-chern-perfect}
to reduce to the case where $E$ is represented
by a bounded complex of finite locally free modules
of fixed rank and apply Lemma \ref{lemma-chern-classes-dual}.
Second proof: use the splitting principle discussed
in Remark \ref{remark-splitting-principle-perfect}
and use that the chern roots of $E^\vee$ are the negatives
of the chern roots of $E$.
\end{proof}

\end{document}
